\chapter{问题模型与理论基础}
\label{ch:theo_base}

\section{探测复杂度与球槽分配模型}

开放寻址哈希表可被抽象为在线状态下的球槽分配问题\cite{aamand2021dynamic,bansal2022deletions,mitzenmacher2001}：表中有固定数量的槽位，每个槽位至多容纳一个元素；系统需要支持动态插入与删除，且任一时刻元素个数不超过槽位数。对任意键 \(x\)，哈希过程会为其生成一条预先确定的随机探测序列
\[
  h_1(x), h_2(x), \ldots
\]
元素只会被安置于该序列给出的候选槽位之一。

当某元素最终被置于其探测序列的第 \(i\) 个候选位置时，探测复杂度（查询时为记录“元素落在第 \(i\) 个候选位置”所需的比特数，记为 \(\Theta(1+\log i)\)，与辅助元数据空间直接相关）定义为 \(\Theta(1+\log i)\)。该量可理解为查询阶段为记录“元素采用了第几个候选位置”而必须保存的比特数目，因而与辅助元数据的空间开销直接相关。在该模型下，用平均探测复杂度刻画整体附加空间，用交换成本刻画单次更新中被搬移元素的个数；最优哈希表理论所关注的核心问题，是在允许有限交换的前提下，如何尽可能降低平均探测复杂度\cite{bender2022otsh,kuszmaul2024filters,pagh2004linear}。一般而言，若表内 \(n\) 个元素分别落在其探测序列的第 \(i_1,i_2,\ldots,i_n\) 个候选位置，则平均探测复杂度为
\[
  \frac{1}{n}\sum_{t=1}^{n}\Theta(1+\log i_t),
\]
而单次插入或删除的交换成本为踢出链中被搬移元素个数的上界。Bender 等人\cite{bender2022otsh} 通过调节 k-kick 深度参数 \(k\)，使上述两项复杂度分别控制在 \(O(\log^{(k)}n)\) 与 \(O(k)\) 量级，形成可调的时间/空间权衡曲线。他们提出：若查询必须在常数时间内完成，且插入、删除允许承担 \(O(k)\) 次元素交换，则每个元素所需的额外空间可由 \(O(\log\log n)\) 比特降至 \(O(\log^{(k)}n)\) 比特。该结果说明，开放寻址哈希表的空间效率不仅取决于空位占比，还与元素搬移、局部引导信息及键恢复编码密切相关。后文将给出与之对应的地址拆分、分层存储与 Local Query Router 设计，使探测序列在插入阶段完整定义、同时满足查询阶段不必线性扫描。

\section{系统核心设定：全域与哈希基础}

\subsection{全局参数}

设键空间全域为 \(U\)，且 \(|U|=\mathrm{poly}(n)\)。当前表内元素规模提示为 \(n\)；表容量 \(N\) 取为 2 的幂，使实际容量落在区间 \([N,2N]\) 内。每个 Facility 的基准规模为
\[
  K=\mathrm{polylog}\,n,
\]
通常取 \(K=\Theta(\log^c n)\) 且 \(c\in[2,4]\)，以便在 Facility 数量 \(N/K\) 与局部结构规模之间取得平衡\cite{bender2022otsh,bender2024sosa}。

\subsection{可逆置换与全局路由哈希}

对任意键 \(x\)，先经随机可逆置换函数 \(\pi\) 映射为 \(\pi(x)\)。该类置换可由 Feistel 等构造从伪随机函数生成\cite{luby1988,carter1979,thorup2012tabulation}。定义 \(x[a, b]\) 表示 \(x\) 的第 \(a\) 位到第 \(b\) 位组成的子串（从低位到高位），则全局路由哈希 \(g^\ast(x)\) 定义为 \(\pi(x)\) 的低 \(\log N\) 位：
\[
  g^\ast(x)=\pi(x)[1,\log N].
\]
置换值在 \(\log N\) 位以上的部分存入槽位载荷，记为
\[
  \mathrm{storage}=\pi(x)[\log N+1,\log U],
\]
即商压缩意义下的 payload；该部分不再参与 Facility 与 Cubby 的路由，仅在命中校验时与局部元数据联立以恢复完整的 \(\pi(x)\)。

从比特布局上看，\(\pi(x)\) 可示意划分为“高位 payload + 低位 \(g^\ast(x)\)”；而 \(g^\ast(x)\) 自身又可自高向低继续划分为 Facility 编号段、Cubby 中间段与 Local Query Router 段，分别用于表级分区、Cubby 内桶化与 \(A[i]\) 索引。图~\ref{fig:gx-layout} 概括了这一层次。

\begin{figure}[htbp]
  \centering
  \setlength{\tabcolsep}{4pt}
  \small
  \begin{tabular}{|c|c|c|} 
    \hline
    \multicolumn{1}{|c|}{$\pi(x)$ 高位：payload / storage} &
    \multicolumn{2}{c|}{$g^\ast(x)=\pi(x)[1,\log N]$} \\
    \hline
    $\pi(x)[\log N+1,\log U]$ & Facility 位 & Cubby 中间位 + $g_I(x)$ \\
    \hline
    存入槽位 storage & $r=g^\ast[\log K+1,\log N]$ & $g^\ast[1,\log K]$ \\
    \hline
  \end{tabular}
  \caption{$\pi(x)$ 与 $g^\ast(x)$ 的位段划分}
  \label{fig:gx-layout}
\end{figure}

\subsection{Facility 编号与 Cubby 内局部地址}

在 \(g^\ast(x)\) 已确定的前提下，Facility 编号取
\[
  r=g^\ast(x)[\log K+1,\log N]
  \;=\;\left\lfloor \frac{g^\ast(x)}{K}\right\rfloor .
\]
对容量为 \(|I|\) 的 Cubby，其内部 Local Query Router 段（与 k-kick 偏好序列无关）定义为
\[
  g_I(x)=g^\ast(x)[1,\log |I|],
\]
即 \(g^\ast(x)\) 的最低 \(\log|I|\) 位，取值范围为 \([0,|I|-1]\)。所有满足 \(g_I(x)=i\) 的键由该 Cubby 内第 \(i\) 号 Local Query Router \(A[i]\) 管理；\(g^\ast(x)\) 中 \([{\log|I|+1},{\log K}]\) 的中间位则与 Facility 编号 \(r\) 及 payload 共同用于恢复完整 \(g^\ast(x)\)，进而经由 \(\pi^{-1}\) 得到原始键 \(x\)\cite{bender2022otsh,pagh2008retrieval}。

\section{核心组件层级关系}

在逻辑结构上，整张表可视为一个大规模数组，均等划分为 \((N/K)\) 个 Facility；每个 Facility 管理规模为 \(K\) 的局部范围，并在其内部维护多层级、不同容量的 Cubby 集合。Cubby 是实际存放元素的局部块；每个 Cubby 内部又包含以下部分：
\begin{itemize}
  \item storage 数组：长度为 \(|I|\) 的槽位数组，存放商压缩 payload 即 \(\pi(x)[\log N+1,\log U]\)；
  \item k-kick 逻辑：将整个 cubby 逻辑上划分为多层级 bin，完成插入、踢出与被踢元素的重插，不要求额外的物理树结构；
  \item MiniArray \(A\)：长度为 \(|I|\) 的 Local Query Router 数组，\(A[i]\) 维护所有 \(g_I(x)=i\) 的键到探测序列下标的映射；
  \item MiniArray \(M\)：记录各 k-kick 深度 bin 的空闲槽位信息，支持在常数时间内判定 bin 内是否存在空位、并定位空位位置；
  \item MiniArray \(B\)：与 storage 槽位一一对应，保存 \(\delta=g_I(x)-i\) 及恢复 \(g^\ast(x)\) 所需的中间位等信息。
\end{itemize}

Facility 层另维护一棵基于 B 树、深度为 \(O(1)\) 的 MiniArray，用于在变长元数据与 Facility 级索引上实现均摊常数时间的访问、插入、删除与更新\cite{bender2022otsh,bayer1972}。MiniArray 中每个内部节点最多有多少个子节点称为 fanout，取 \(\Theta(\log n/\log\log n)\)，内部指针仅需 \(\Theta(\log\log n)\) 比特，从而到达整体 \(O(\log^{(k)}n)\) 的空间目标\cite{pagh2001,jacobson1989}。

为在动态扩缩容下维持高装填因子，每个 Facility 还引入 \emph{tail Cubby} 机制：初始时仅含一个 1-tiered Cubby 作为 tail；此后所有新插入默认进入 tail，tail 满后转为普通 1-tiered Cubby 并新建 tail。不同 \(j\)-tiered Cubby 的目标个数与容量分别为
\[
  t_j=\Theta\!\left(\frac{(\log^{(j+1)}n)^2}{(\log^{(j)}n)^2}\right),\qquad
  r_j=\frac{K}{(\log^{(j)}n)^2};
\]
当某层 Cubby 数量偏离 \(t_j\) 时，通过向上合并或向下拆分并在新 Cubby 内重插全部元素来恢复分布不变量。删除非 tail 元素时，从 tail 随机抽取一个元素回填至被删槽位；tail 为空时从 1-tiered 层提升一个 Cubby 作为新 tail，并可能触发拆分。tier 调度、tail 不变量与 k-kick 踢出规则共同构成第三章的算法主体\cite{bender2022otsh,bender2024sosa,conway2018,pandey2023iceberght}。

\begin{figure}[htbp]
  \centering
  \begin{tikzpicture}[
    node distance=0.65cm,
    box/.style={draw, rounded corners=2pt, minimum height=0.7cm, align=center, font=\small, inner sep=3pt},
    top/.style={box},
    child/.style={box, minimum width=2.55cm, minimum height=0.82cm},
    arr/.style={-{Stealth[length=2mm]}, thick}
  ]
    \node[top, minimum width=3.8cm] (ht) {哈希表};
    \node[top, minimum width=6cm, below=of ht] (fac) {Facility（$N/K$ 个表级分区）};
    \node[top, minimum width=6cm, below=of fac] (cub) {Cubby（多层级 + tail）};
    \node[child, anchor=north] (kick) at ([xshift=-4.5cm, yshift=-1.55cm]cub.south) {k-kick\\逻辑 bin};
    \node[child, anchor=north] (router) at ([xshift=-1.5cm, yshift=-1.55cm]cub.south) {Local Query Router\\（$A[i]$ 前缀树）};
    \node[child, anchor=north] (M) at ([xshift=1.5cm, yshift=-1.55cm]cub.south) {MiniArray $M$\\（空闲位图）};
    \node[child, anchor=north] (B) at ([xshift=4.5cm, yshift=-1.55cm]cub.south) {MiniArray $B$\\（键恢复差分）};
    \draw[arr] (ht) -- (fac);
    \draw[arr] (fac) -- (cub);
    \coordinate (bus) at ([yshift=-0.55cm]cub.south);
    \draw[arr] (cub.south) -- (bus);
    \draw[arr] (bus) -| (kick.north);
    \draw[arr] (bus) -| (router.north);
    \draw[arr] (bus) -| (M.north);
    \draw[arr] (bus) -| (B.north);
  \end{tikzpicture}
  \caption{Facility--Cubby--Slot 分层架构}
  \label{fig:sys-arch}
\end{figure}

\section{商压缩的信息论思想}

如果在槽位中完整保存原始键，则部分信息与寻址过程重复。商压缩借助双射函数 \(\pi\)，映射 \(x\) 为 \(\pi(x)\)，低 \(\log N\) 位 \(\pi(x)[1,\log N]=g^\ast(x)\) 的各位段交由地址上下文（Facility 编号、Cubby 容量、槽位下标）与局部元数据分担，槽位仅保存 payload 即 \(\pi(x)[\log N+1,\log U]\) 与无法由上下文恢复的差分信息。

具体地，当 storage\([i]\) 存放键 \(x\) 对应信息时，MiniArray \(B[i]\) 保存
\[
  \delta=g_I(x)-i,
\]
以及 \(g^\ast(x)[\log|I|+1,\log K]\) 的比特位。配合 Facility 编号 \(r=g^\ast(x)[\log K+1,\log N]\) ，即可重建
\[
  g^\ast(x)[1,\log N],
\]
再加上 payload 得到 \(\pi(x)\)，经 \(\pi^{-1}\) 恢复 \(x\)。该编码方式与紧凑检索结构\cite{pagh2008retrieval,lehmann2026mphsurvey}、打包式哈希表\cite{kurpicz2023pachash,hermann2024phobic} 所强调的“按条目实际长度计费”相同；但在 k-kick 动态搬移与 tier 重建下，\(B\)、\(A\) 与 \(M\) 必须同步更新，对 MiniArray 与 Local Query Router 的动态维护提出了比静态紧凑哈希更严格的要求\cite{bender2022otsh,kuszmaul2024filters}。

\section{k-kick、MiniArray 与 Local Query Router}

\subsection{k-kick：逻辑分层与交换上界}

在 Cubby 内部，k-kick 并不引入物理上的多层数组，而是将 storage 上长度为 \(|I|\) 的连续槽位按固定规则划分为 \(k+1\) 个逻辑深度：深度 0 的 bin （ bin 指 Cubby 按不同深度划分的子区间）覆盖整个 Cubby，更深层 bin 由上一层均匀切分得到，规模 \(s_i=\Theta((\log^{(i)}K)^6)\)。对键 \(x\)（此处均指 \(\pi(x)\)）定义偏好序列
\[
  g_0(x)\rightarrow g_1(x)\rightarrow\cdots\rightarrow g_k(x),
\]
其中 \(g_0(x)\) 为整个 Cubby，\(g_{i+1}(x)\) 为 \(g_i(x)\) 经均匀切分后 \(x\) 所属的子 bin。探测位置序列 \(h_1(x),h_2(x),\ldots\) 按先深后浅顺序枚举各 bin 内槽位；若元素最终落在第 \(j\) 个候选位置，则其探测复杂度为 \(\Theta(1+\log j)\)。

定义 bin 在已满且其中每个元素的 insert\_depth 均不低于该 bin 对应深度时饱和（IsSaturated）；贪心策略下 GetMaxUsableDepth(\(x\)) 返回 \(x\) 可插入的、未饱和 bin 的最大深度。插入时，系统先取随机哈希 \(s(x)\in[0,k]\)，再令
\[
  \mathrm{depth}=\min\bigl(\Call{GetMaxUsableDepth}{x},\,s(x)\bigr),
\]
在 bin \(g_{\mathrm{depth}}(x)\) 内写入；若 bin 已满，则踢出其中 insert\_depth 最小的元素并调用 ReInsert 向上安置。Bender 等人\cite{bender2022otsh} 证明：在参数满足其设定时，整条踢出链至多 \(k\) 次，单次插入的交换成本为 \(O(k)\)；配合随机深度选择，平均探测复杂度可控制在 \(O(\log^{(k)}n)\) 量级，从而与 上述描述的空间—时间权衡曲线对齐\cite{pagh2004cuckoo,lehmann2023sichash,pagh2004linear}。

\subsection{MiniArray \(M\)}

k-kick 的 IsSaturated、GetMaxUsableDepth 与 ReInsert 都需要在 bin 上判断是否存在空位并找到可用位置。MiniArray \(M\) 为各深度 bin 维护等价于 bitmap 的空闲槽位图\cite{wang2023bitmap}，使上述判定与定位在均摊意义下为常数时间，而不必扫描整个 Cubby\cite{bender2022otsh,pagh2004linear}。 MiniArray 则通过 \(O(1)\) 深度 B 树、叶内 Rank/Select 与 Packed Leaf 维护变长元数据，并实现常数时间复杂度的 Access/Insert/Delete/Update 函数\cite{bayer1972,jacobson1989,pagh2001}：树高 \(O(1)\)，内部指针 \(\Theta(\log\log n)\) 比特\cite{kurpicz2023pachash,lehmann2024shockhash}。

\subsection{Local Query Router：\(O(1)\) 查询与映射存储}

偏好序列与探测位置序列在插入阶段完整定义；为实现 $O(1)$ 查询，须额外维护键到探测下标的映射，而不遍历 $h_1(x), h_2(x), \dots$：k-kick 插入会搬移元素，其实际位置既不等于首选位置 $g_k(x)$，也不等于 $h_1(x)$。Local Query Router 因此为每个键保存映射
\[
  \pi(x)\longrightarrow j,
\]
其中 \(j\) 为探测序列下标。Cubby 内令 \(g_I(x)=g^\ast(x)[1,\log|I|]\)，所有满足 \(g_I(x)=i\) 的键由 router \(A[i]\) 管理；\(A[i]\) 实现为二叉前缀树（可进一步 Patricia 压缩\cite{morrison1968,lehmann2025combined}），对外接口为 insert(key, j)、query(key)、delete(key)。查询时先算 \(i=g_I(x)\)，在 \(A[i]\) 得 \(j\)，再访问 storage\([h_j(x)]\)，加上 \(B[h_j(x)]\) 与 payload 恢复 \(\pi(x)\)，与查询侧的 \(\pi(x)\) 进行对比；\(j\le\log K\)，每条映射仅需 \(O(\log\log n)\) 比特，查询路径为常数次 router 与 MiniArray 访问\cite{bender2022otsh,pagh2008retrieval}。

\subsection{辅助结构同步不变量}

元素在 k-kick 踢出链中每发生一次搬移，须同步更新四处状态：
\begin{itemize}
  \item storage 槽位内容与 insert\_depth；
  \item \(B[i]\) 中的 \(\delta\) 与中间位；
  \item \(A[g_I(x)]\) 中 \(\pi(x)\mapsto j\) 的映射；
  \item \(M\) 中各 bin 的空闲位图。
\end{itemize}
删除非 tail 元素时的 tail 随机回填、tier 合并/拆分后的全量重插，以及表级双表扩缩容，均遵循同一同步规则——不可以仅复制 payload 而忽略 router 与 \(B\) 的上下文依赖\cite{bender2022otsh,bender2024sosa}。第三章将给出 MiniArray、k-kick 与 Local Query Router 的全部算法伪代码及统一操作语义。

\section{本章小结}

本章建立了后文所需的理论语境：球槽分配模型给出探测复杂度与交换成本的定义\cite{knuth1973sorting,aamand2021dynamic}；全域参数与 \(g^\ast(x)\) 位段划分明确了地址语义；图~\ref{fig:sys-arch} 与 Facility--Cubby--Slot(\(storage/A/M/B\)) 层级关系给出整体骨架；商压缩说明 payload 与 \(B[i]\) 的信息分工\cite{bender2022otsh,lehmann2025combined}；同时从 k-kick 交换上界、MiniArray \(M\) 的空闲判定、Local Query Router 的 \(O(1)\) 查询及四个状态同步不变量四个角度，为第三章算法设计提供理论锚点。下一章将给出各模块的完整伪代码；第四章则报告宏基准与参数扫描下的性能测量。
